\answer{ex:12:1}{$0.60$ (sem correlação seria $0.18$), limite $\sqrt{c/(rT)}\approx0.58$: cem neurônios distinguem apenas ``sim/não''.}
\answer{ex:12:2}{(a)~$10^8$ disparos, $\approx10\,\text{mJ}$: $0.3\,\%$ dos $3$~J do cérebro inteiro. (b)~$3$~J, $\approx300$ vezes mais.}
\answer{ex:12:3}{(a)~$\theta_A\in[2,3)$, $\theta_B\in[1,2)$; a figura alheia dá no máximo $2$ e $1$. (b)~$10\cdot96^2\approx9.2\cdot10^4$.}
\answer{ex:12:4}{(a)~$T_d=0.85\,\tau_a=0.34\,\text{s}$. (b)~$T_d\propto\tau_a$ e não depende de $I$: $\tau_a\approx3.5\,\text{s}$ (simulação: $3.1\,\text{s}$).}
\answer{ex:12:5}{$p\approx0.17$. A fração cai devagar: $0.90$, $0.88$, $0.86$, $0.84$, $0.82$, $0.79$ com $N=8$, $12$, $24$, $48$, $96$, $192$.}
\answer{ex:12:6}{As dez execuções não têm erro com $\tau_{\text{tr}}=20$, $50$ e $200\ms$ (com $30\ms$, uma execução em dez falhou); com $5$--$10\ms$, cerca de $0.5$; com $0.5$--$2\,\text{s}$, a qualidade cai para $0.86$. Por cima: o traço deve se apagar durante a pausa, $e^{-0.3\,\text{s}/\tau_{\text{tr}}}\ll1$.}
\answer{ex:12:7}{Um só atraso não basta: a janela de $\pm5\ms$ não cobre $\Delta x/v$ de $5$ a $20\ms$. Dois ramos de inibição com $5<d_1\le10\ms$ e $15\le d_2\le d_1+10\ms$.}
\answer{ex:12:8}{\texttt{two\_stage}: $\Delta=2$, empate com uma inversão, troca de resposta a partir de duas. Contador de uns: uma inversão; $0.57$ com $p=0.1$.}
